\chapter{Nombres complexes}

\noindent\textit{Inventés pour résoudre les équations que les réels ne savaient pas traiter
(comme $x^{2}=-1$), les nombres complexes offrent un pont remarquable entre l'algèbre et la
géométrie : tout point du plan devient un nombre, et toute transformation géométrique une
opération sur ces nombres. C'est l'un des chapitres emblématiques de la série C, où l'on calcule
des modules et des arguments, où l'on linéarise des puissances de cosinus, et où l'on démontre
qu'un triangle est rectangle isocèle en une seule ligne de calcul.}
\vspace{8pt}

% ═══════════════════════════════════════════════════════════════
\section{Forme algébrique}

\begin{definition}[L'ensemble $\C$]
Il existe un ensemble $\C$, contenant $\R$, muni d'une addition et d'une multiplication qui
prolongent celles de $\R$, et d'un élément $i$ vérifiant $i^{2}=-1$, tel que tout élément $z$ de
$\C$ s'écrit de façon \textbf{unique}
\[
z=a+ib,\qquad (a,b)\in\R^{2}.
\]
$a=\mathrm{Re}(z)$ est la \textbf{partie réelle}, $b=\mathrm{Im}(z)$ la \textbf{partie
imaginaire} (toutes deux \emph{réelles}). Si $b=0$, $z$ est \textbf{réel} ; si $a=0$, $z$ est un
\textbf{imaginaire pur}.
\end{definition}

\begin{remarque}
L'écriture étant unique, deux complexes sont égaux si et seulement si ils ont \emph{même partie
réelle et même partie imaginaire} :
\[
a+ib=a'+ib'\iff a=a'\ \text{et}\ b=b'.
\]
C'est le principe d'\textbf{identification}, fondamental dans la résolution d'équations.
\end{remarque}

\begin{propriete}[Opérations sous forme algébrique]
Pour $z=a+ib$ et $z'=a'+ib'$ :
\begin{itemize}
\item $z+z'=(a+a')+i(b+b')$ ;
\item $z\,z'=(aa'-bb')+i(ab'+a'b)$ (on développe et on remplace $i^{2}$ par $-1$) ;
\item si $z\neq0$, $\dfrac1z=\dfrac{a-ib}{a^{2}+b^{2}}$ (on multiplie par le conjugué).
\end{itemize}
\end{propriete}

\begin{exemple}
$(2-i)(3+i)=6+2i-3i-i^{2}=6-i+1=7-i$. Et
$\dfrac{1}{2-i}=\dfrac{2+i}{(2-i)(2+i)}=\dfrac{2+i}{5}=\dfrac25+\dfrac15 i$.
\end{exemple}

\begin{aretenir}
Les puissances de $i$ sont \textbf{périodiques de période $4$} :
\[
i^{0}=1,\quad i^{1}=i,\quad i^{2}=-1,\quad i^{3}=-i,\quad i^{4}=1,\ \dots
\]
Pour calculer $i^{n}$, il suffit de regarder le reste de $n$ dans la division par $4$. Par exemple
$i^{2027}=i^{4\times506+3}=i^{3}=-i$.
\end{aretenir}

\subsection{Conjugué}

\begin{definition}[Conjugué]
Le \textbf{conjugué} de $z=a+ib$ est $\bar z=a-ib$. Géométriquement, $\bar z$ est le symétrique de
$z$ par rapport à l'axe des réels.
\end{definition}

\begin{propriete}[Règles de conjugaison]
$\overline{z+z'}=\bar z+\bar z'$,\quad $\overline{z\,z'}=\bar z\,\bar z'$,\quad
$\overline{\left(\dfrac1z\right)}=\dfrac1{\bar z}$,\quad $\overline{z^{n}}=(\bar z)^{n}$,\quad
$\overline{\bar z}=z$.\\[2pt]
De plus :
\[
z+\bar z=2\,\mathrm{Re}(z),\qquad z-\bar z=2i\,\mathrm{Im}(z),\qquad z\,\bar z=a^{2}+b^{2}.
\]
En particulier : $z\in\R\iff z=\bar z$, et $z$ imaginaire pur $\iff z=-\bar z$.
\end{propriete}

\begin{remarque}
La relation $z\,\bar z=a^{2}+b^{2}$ est le moteur de toute division : pour rendre réel un
dénominateur, on le multiplie (haut et bas) par son conjugué. C'est l'analogue de la
« quantité conjuguée » utilisée avec les racines carrées.
\end{remarque}

\subsection{Module}

\begin{definition}[Module]
Le \textbf{module} de $z=a+ib$ est le réel positif
\[
|z|=\sqrt{a^{2}+b^{2}}=\sqrt{z\,\bar z}.
\]
C'est la distance de l'origine au point d'affixe $z$. On a $|z|=0\iff z=0$.
\end{definition}

\begin{illus}
\begin{tikzpicture}[>=Stealth,scale=1.1]
\draw[->] (-0.4,0) -- (4,0) node[right]{\small $\mathrm{Re}$};
\draw[->] (0,-0.4) -- (0,2.8) node[above]{\small $\mathrm{Im}$};
\filldraw[onbo] (3,2) circle (2pt) node[above right]{\small $M(z)$};
\draw[onbo,->,line width=1pt] (0,0) -- (3,2);
\draw[dashed,onbmuted] (3,0) node[below]{\small $a$} -- (3,2);
\draw[dashed,onbmuted] (0,2) node[left]{\small $b$} -- (3,2);
\filldraw[onbmuted] (3,-2) circle (1.4pt) node[below right]{\small $\bar z$};
\draw[onbmuted,dotted] (3,2)--(3,-2);
\node[onbgreen] at (1.3,1.25) {\small $|z|$};
\draw[onbgreen] (0.8,0) arc (0:33.7:0.8);
\node[onbgreen] at (1.05,0.22) {\small $\theta$};
\node[align=left,text width=4.4cm] at (6.4,1.1)
{\small À $z=a+ib$ on associe le point $M(a,b)$ d'\textbf{affixe} $z$.
Alors $|z|=OM$, $\theta=\arg(z)=(\vec{u},\overrightarrow{OM})$, et $\bar z$ est le symétrique de
$z$ par rapport à l'axe réel.};
\end{tikzpicture}
\fig{Figure — Le plan complexe : affixe, module, argument et conjugué.}
\end{illus}

\begin{propriete}[Propriétés du module]
Pour tous $z,z'\in\C$ :
\[
|z\,z'|=|z|\,|z'|,\quad \left|\dfrac z{z'}\right|=\dfrac{|z|}{|z'|}\ (z'\neq0),\quad
|z^{n}|=|z|^{n},\quad |\bar z|=|z|.
\]
\textbf{Inégalité triangulaire} : $\,|z+z'|\le|z|+|z'|$, avec égalité si et seulement si $z$ et
$z'$ sont « de même direction » ($z'=\lambda z$, $\lambda\ge0$).
\end{propriete}

\begin{theoreme}[Équation du second degré à coefficients réels]
L'équation $az^{2}+bz+c=0$ ($a\neq0$, $a,b,c$ réels) de discriminant $\Delta=b^{2}-4ac$ :
\begin{itemize}
\item si $\Delta\ge0$ : deux racines réelles (ou une double) usuelles ;
\item si $\Delta<0$ : deux solutions complexes \emph{conjuguées}
\[
z_{1,2}=\frac{-b\pm i\sqrt{-\Delta}}{2a}.
\]
\end{itemize}
\end{theoreme}

\begin{exemple}
$z^{2}-2z+5=0$ : $\Delta=4-20=-16<0$, $\sqrt{-\Delta}=4$, donc $z=\dfrac{2\pm 4i}{2}=1\pm2i$. Les
deux racines sont bien conjuguées.
\end{exemple}

\begin{remarque}
Lorsqu'un polynôme à coefficients réels admet une racine complexe $z_0$, il admet aussi $\bar z_0$.
On peut alors factoriser par $\big(z-z_0\big)\big(z-\bar z_0\big)=z^{2}-2\,\mathrm{Re}(z_0)\,z+|z_0|^{2}$.
\end{remarque}

% ═══════════════════════════════════════════════════════════════
\section{Formes trigonométrique et exponentielle}

\begin{definition}[Argument, forme trigonométrique]
Pour $z\neq0$ de module $r=|z|$, un \textbf{argument} $\theta=\arg(z)$ (défini \emph{modulo}
$2\pi$) est un réel vérifiant
\[
\cos\theta=\dfrac ar,\qquad \sin\theta=\dfrac br.
\]
On obtient la \textbf{forme trigonométrique}
\[
z=r\big(\cos\theta+i\sin\theta\big).
\]
\end{definition}

\begin{definition}[Notation exponentielle]
On pose, pour tout $\theta\in\R$,
\[
e^{i\theta}=\cos\theta+i\sin\theta.
\]
Tout complexe non nul s'écrit alors sous \textbf{forme exponentielle} $z=r\,e^{i\theta}$, avec
$r=|z|>0$ et $\theta=\arg(z)$. On a $|e^{i\theta}|=1$ et $\overline{e^{i\theta}}=e^{-i\theta}$.
\end{definition}

\begin{theoreme}[Relation fonctionnelle de l'exponentielle complexe]
Pour tous $\theta,\theta'\in\R$ :
\[
e^{i\theta}\,e^{i\theta'}=e^{i(\theta+\theta')},\qquad \frac{1}{e^{i\theta}}=e^{-i\theta},\qquad
\big(e^{i\theta}\big)^{n}=e^{in\theta}\ (n\in\Z).
\]
\end{theoreme}

\noindent\emph{Justification (produit).} En développant et en utilisant les formules d'addition :
\begin{align*}
e^{i\theta}e^{i\theta'}&=(\cos\theta+i\sin\theta)(\cos\theta'+i\sin\theta')\\
&=(\cos\theta\cos\theta'-\sin\theta\sin\theta')+i(\sin\theta\cos\theta'+\cos\theta\sin\theta')\\
&=\cos(\theta+\theta')+i\sin(\theta+\theta')=e^{i(\theta+\theta')}.
\end{align*}

\begin{propriete}[Module et argument d'un produit, d'un quotient]
Pour $z,z'$ non nuls (égalités sur les arguments \emph{modulo} $2\pi$) :
\[
\arg(z\,z')=\arg z+\arg z',\qquad \arg\!\dfrac z{z'}=\arg z-\arg z',
\]
\[
\arg(\bar z)=-\arg z,\qquad \arg(z^{n})=n\arg z.
\]
Multiplier par $z'$ revient donc à \textbf{multiplier les modules} et \textbf{ajouter les
arguments} : c'est une similitude (rotation + homothétie).
\end{propriete}

\begin{illus}
\begin{tikzpicture}[>=Stealth,scale=1.25]
\draw[onbline] (0,0) circle (1);
\draw[->] (-1.4,0)--(1.5,0) node[right]{\small $\mathrm{Re}$};
\draw[->] (0,-1.0)--(0,1.5) node[above]{\small $\mathrm{Im}$};
\filldraw[onbo] (40:1) circle (1.3pt);
\draw[onbo,->] (0,0)--(40:1) node[midway,below right]{\scriptsize $z$};
\filldraw[onbblue] (100:1) circle (1.3pt);
\draw[onbblue,->] (0,0)--(100:1) node[midway,above left]{\scriptsize $z\cdot z'$};
\draw[onbgreen] (40:0.45) arc (40:100:0.45);
\node[onbgreen] at (70:0.65) {\scriptsize $+\arg z'$};
\node[align=left,text width=4.0cm] at (3.4,0.2)
{\small Sur le cercle unité, multiplier par $z'=e^{i\alpha}$ \textbf{fait tourner} d'un angle
$\alpha=\arg z'$ : les arguments s'ajoutent.};
\end{tikzpicture}
\fig{Figure — Multiplier par $e^{i\alpha}$ est une rotation d'angle $\alpha$.}
\end{illus}

\begin{theoreme}[Formule de Moivre]
Pour tout $\theta\in\R$ et tout $n\in\Z$ :
\[
\big(\cos\theta+i\sin\theta\big)^{n}=\cos(n\theta)+i\sin(n\theta).
\]
\end{theoreme}

\noindent\emph{Justification.} En notation exponentielle, c'est simplement
$\big(e^{i\theta}\big)^{n}=e^{in\theta}$, ce qui donne, en repassant en forme trigonométrique,
$\cos(n\theta)+i\sin(n\theta)$.

\begin{theoreme}[Formules d'Euler]
Pour tout $\theta\in\R$ :
\[
\cos\theta=\frac{e^{i\theta}+e^{-i\theta}}{2},\qquad
\sin\theta=\frac{e^{i\theta}-e^{-i\theta}}{2i}.
\]
\end{theoreme}

\noindent\emph{Justification.} On additionne (resp. soustrait) $e^{i\theta}=\cos\theta+i\sin\theta$
et $e^{-i\theta}=\cos\theta-i\sin\theta$.

\begin{aretenir}
Les formes $e^{i\theta}$ sont idéales pour les \textbf{puissances} (Moivre) ; les formules d'Euler
servent à \textbf{linéariser} ($\cos^{n}\theta$, $\sin^{n}\theta$, $\cos^{p}\theta\sin^{q}\theta$)
et à transformer des produits en sommes. Pensez : « Moivre monte en puissance, Euler descend en
somme ».
\end{aretenir}

\begin{exemple}[Valeurs remarquables]
\[
e^{i0}=1,\quad e^{i\pi/2}=i,\quad e^{i\pi}=-1,\quad e^{-i\pi/2}=-i,\quad e^{2i\pi}=1.
\]
L'identité $e^{i\pi}+1=0$ relie cinq constantes fondamentales des mathématiques.
\end{exemple}

% ═══════════════════════════════════════════════════════════════
\section{Racines $n$-ièmes}

\begin{theoreme}[Racines $n$-ièmes de l'unité]
Soit $n\ge1$ entier. L'équation $z^{n}=1$ admet exactement $n$ solutions, les \textbf{racines
$n$-ièmes de l'unité} :
\[
\omega_{k}=e^{\,2ik\pi/n},\qquad k\in\{0,1,\dots,n-1\}.
\]
Leurs images sont les sommets d'un \textbf{polygone régulier} à $n$ côtés inscrit dans le cercle
unité, l'un d'eux étant le point $1$. Leur \textbf{somme} vaut $0$ (pour $n\ge2$).
\end{theoreme}

\begin{theoreme}[Racines $n$-ièmes d'un complexe non nul]
Soit $Z=R\,e^{i\varphi}$ avec $R>0$. L'équation $z^{n}=Z$ admet exactement $n$ solutions :
\[
z_{k}=R^{1/n}\,e^{\,i(\varphi+2k\pi)/n},\qquad k\in\{0,1,\dots,n-1\}.
\]
On obtient toutes les solutions en multipliant l'une d'elles par les racines $n$-ièmes de l'unité.
\end{theoreme}

\begin{illus}
\begin{tikzpicture}[>=Stealth,scale=1.3]
\draw[onbline] (0,0) circle (1);
\draw[->] (-1.3,0)--(1.4,0) node[right]{\small $\mathrm{Re}$};
\draw[->] (0,-1.3)--(0,1.4) node[above]{\small $\mathrm{Im}$};
\foreach \k in {0,...,5}{
  \filldraw[onbo] (\k*60:1) circle (1.4pt);
}
\draw[onbo,line width=0.7pt] (0:1)--(60:1)--(120:1)--(180:1)--(240:1)--(300:1)--cycle;
\node[onbo,right] at (0:1) {\scriptsize $\omega_0{=}1$};
\node[onbo,above right] at (60:1) {\scriptsize $\omega_1$};
\node[align=left,text width=4.1cm] at (4.0,0)
{\small Les $6$ racines sixièmes de l'unité $e^{ik\pi/3}$ forment un hexagone régulier inscrit
dans le cercle unité ; leur somme est nulle.};
\end{tikzpicture}
\fig{Figure — Racines sixièmes de l'unité : un polygone régulier.}
\end{illus}

\begin{definition}[Racine carrée d'un complexe sous forme algébrique]
Pour résoudre $z^{2}=Z$ avec $Z=\alpha+i\beta$ \emph{sans} forme trigonométrique, on pose
$z=x+iy$ et on résout le système
\[
x^{2}-y^{2}=\alpha,\qquad 2xy=\beta,\qquad x^{2}+y^{2}=|Z|=\sqrt{\alpha^{2}+\beta^{2}},
\]
(la troisième équation étant $|z|^{2}=|Z|$). Le signe de $\beta$ fixe le signe relatif de $x$ et
$y$. Cette méthode est précieuse dans les équations du second degré à coefficients complexes.
\end{definition}

\begin{exemple}
Racines carrées de $Z=3+4i$. On a $|Z|=5$, donc $x^{2}-y^{2}=3$, $x^{2}+y^{2}=5$, d'où $x^{2}=4$,
$y^{2}=1$. Comme $2xy=4>0$, $x$ et $y$ sont de même signe : $z=\pm(2+i)$.
\end{exemple}

% ═══════════════════════════════════════════════════════════════
\section{Équations dans $\C$}

\begin{theoreme}[Équation du second degré à coefficients complexes]
L'équation $az^{2}+bz+c=0$ ($a\neq0$, $a,b,c\in\C$) se résout par
\[
z=\frac{-b\pm\delta}{2a},
\]
où $\delta$ est \emph{une} racine carrée complexe de $\Delta=b^{2}-4ac$ (déterminée par la méthode
algébrique ci-dessus). Le « $\pm$ » garde tout son sens car $\delta$ et $-\delta$ sont les deux
racines carrées de $\Delta$.
\end{theoreme}

\begin{remarque}
Le théorème de d'Alembert--Gauss (admis) affirme que tout polynôme de degré $n\ge1$ à coefficients
complexes possède exactement $n$ racines dans $\C$ (comptées avec multiplicité). $\C$ est dit
\textbf{algébriquement clos} : toute équation polynomiale y a des solutions.
\end{remarque}

\begin{exemple}
$z^{2}-(3+i)z+(2+2i)=0$. $\Delta=(3+i)^{2}-4(2+2i)=9+6i-1-8-8i=-2i$. Une racine carrée de $-2i$ :
on cherche $\delta=x+iy$ avec $x^{2}-y^{2}=0$, $2xy=-2$, $x^{2}+y^{2}=2$, d'où $x^{2}=y^{2}=1$ et
$xy<0$ : $\delta=1-i$. Donc $z=\dfrac{(3+i)\pm(1-i)}{2}$, soit $z=2$ ou $z=1+i$.
\end{exemple}

% ═══════════════════════════════════════════════════════════════
\section{Interprétation géométrique}

\begin{definition}[Affixe d'un point, d'un vecteur]
Dans le plan muni d'un repère orthonormé direct $(O,\vec u,\vec v)$, le point $M(a,b)$ a pour
\textbf{affixe} $z_M=a+ib$. Le vecteur $\vec w(a,b)$ a pour affixe $a+ib$. L'affixe d'une somme de
vecteurs est la somme des affixes ; l'affixe du milieu $I$ de $[AB]$ est
$z_I=\dfrac{z_A+z_B}{2}$.
\end{definition}

\begin{propriete}[Distances et angles]
Pour $A,B,C$ d'affixes $z_A,z_B,z_C$ :
\[
\overrightarrow{AB}\ \text{a pour affixe}\ z_B-z_A,\qquad AB=|z_B-z_A|,\qquad
(\vec u,\overrightarrow{AB})=\arg(z_B-z_A).
\]
De plus
\[
\frac{z_C-z_A}{z_B-z_A}\ \text{a pour module}\ \frac{AC}{AB}\ \text{et pour argument}\
(\overrightarrow{AB},\overrightarrow{AC}).
\]
\end{propriete}

\begin{aretenir}
Le rapport $\dfrac{z_C-z_A}{z_B-z_A}$ est l'outil clé des configurations :
\begin{itemize}
\item \textbf{réel} $\iff$ $A,B,C$ \textbf{alignés} ;
\item \textbf{imaginaire pur} $\iff$ $(AB)\perp(AC)$ ;
\item de \textbf{module $1$} $\iff$ $AB=AC$ (triangle isocèle en $A$) ;
\item égal à $i$ $\iff$ triangle \textbf{rectangle isocèle} direct en $A$.
\end{itemize}
\end{aretenir}

\begin{illus}
\begin{tikzpicture}[>=Stealth,scale=1.2]
\coordinate (A) at (0,0);
\coordinate (B) at (2.4,0);
\coordinate (C) at (0,2.4);
\draw[->] (-0.5,0)--(3,0) node[right]{\small $\mathrm{Re}$};
\draw[->] (0,-0.5)--(0,3) node[above]{\small $\mathrm{Im}$};
\filldraw[onbo] (A) circle (1.4pt) node[below left]{\small $A$};
\filldraw[onbo] (B) circle (1.4pt) node[below right]{\small $B$};
\filldraw[onbo] (C) circle (1.4pt) node[above left]{\small $C$};
\draw[onbo,line width=0.8pt] (B)--(A)--(C);
\draw[onbgreen] (0.45,0) arc (0:90:0.45);
\node[onbgreen] at (0.7,0.7) {\small $\frac\pi2$};
\node[align=left,text width=4.0cm] at (3.4,1.4)
{\small Ici $\dfrac{z_C-z_A}{z_B-z_A}=\dfrac{2{,}4\,i}{2{,}4}=i$ : argument $\frac\pi2$ (angle
droit) et module $1$ ($AB=AC$). Le triangle est rectangle isocèle en $A$.};
\end{tikzpicture}
\fig{Figure — Lecture d'une configuration par un quotient d'affixes.}
\end{illus}

\begin{propriete}[Ensembles de points classiques]
\begin{itemize}
\item $|z-z_A|=r$ : \textbf{cercle} de centre $A$, rayon $r$ ;
\item $|z-z_A|=|z-z_B|$ : \textbf{médiatrice} de $[AB]$ ;
\item $\arg(z-z_A)=\theta_0\ [2\pi]$ : \textbf{demi-droite} d'origine $A$ ;
\item $\dfrac{z-z_A}{z-z_B}\in\R$ : la \textbf{droite} $(AB)$ (privée de $B$).
\end{itemize}
\end{propriete}

% ═══════════════════════════════════════════════════════════════
\section{L'essentiel à retenir}

\begin{synthese}
\textbf{Forme algébrique.} $z=a+ib$, unique. $\mathrm{Re}(z)=a$, $\mathrm{Im}(z)=b$. Égalité
$\iff$ mêmes parties réelle et imaginaire. $i^{2}=-1$, puissances de $i$ de période $4$.

\smallskip
\textbf{Conjugué \& module.} $\bar z=a-ib$ ;\ $z\bar z=|z|^{2}=a^{2}+b^{2}$ ;\
$z+\bar z=2\mathrm{Re}(z)$,\ $z-\bar z=2i\,\mathrm{Im}(z)$. \\
$|zz'|=|z||z'|$,\ $|z/z'|=|z|/|z'|$,\ $|z^{n}|=|z|^{n}$,\ $|z+z'|\le|z|+|z'|$.

\smallskip
\textbf{Formes trigo./exp.} $z=r(\cos\theta+i\sin\theta)=re^{i\theta}$, $r=|z|$,
$\theta=\arg z\ [2\pi]$. \\
$\arg(zz')=\arg z+\arg z'$,\ $\arg(z/z')=\arg z-\arg z'$,\ $\arg(z^{n})=n\arg z$,\
$\arg\bar z=-\arg z$.

\smallskip
\textbf{Moivre \& Euler.} $(\cos\theta+i\sin\theta)^{n}=\cos n\theta+i\sin n\theta$ ;\
$\cos\theta=\dfrac{e^{i\theta}+e^{-i\theta}}2$,\ $\sin\theta=\dfrac{e^{i\theta}-e^{-i\theta}}{2i}$.
Moivre $\to$ puissances ; Euler $\to$ linéarisation.

\smallskip
\textbf{Racines $n$-ièmes.} $z^{n}=Re^{i\varphi}\Rightarrow z_k=R^{1/n}e^{i(\varphi+2k\pi)/n}$,
$k=0,\dots,n-1$ ($n$ solutions, polygone régulier ; somme nulle pour l'unité).

\smallskip
\textbf{Second degré.} Réel $\Delta<0$ : $z=\dfrac{-b\pm i\sqrt{-\Delta}}{2a}$. Complexe : $\delta$
racine carrée de $\Delta$ par identification, $z=\dfrac{-b\pm\delta}{2a}$.

\smallskip
\textbf{Géométrie.} $z_B-z_A\leftrightarrow\overrightarrow{AB}$ ; $AB=|z_B-z_A|$ ;
$(\overrightarrow{AB},\overrightarrow{AC})=\arg\dfrac{z_C-z_A}{z_B-z_A}$. Quotient réel $=$
alignés ; imaginaire pur $=$ orthogonaux ; module $1=$ isocèle en $A$.

\smallskip
\textbf{Pièges.} L'argument n'est défini qu'à $2\pi$ près (et $0$ n'en a pas). Vérifier le signe de
$\sin\theta$ \emph{et} de $\cos\theta$ pour placer $\theta$. Ne pas écrire $\sqrt{\Delta}$ quand
$\Delta\in\C$ : il y a \emph{deux} racines carrées, sans relation d'ordre.
\end{synthese}

% ═══════════════════════════════════════════════════════════════
\section{Méthodes \& savoir-faire}

\methode{Mettre un quotient sous forme algébrique}
Multiplier numérateur et dénominateur par le \textbf{conjugué du dénominateur}, puis simplifier.
\begin{exemple}
$\dfrac{3+i}{1-i}=\dfrac{(3+i)(1+i)}{(1-i)(1+i)}=\dfrac{3+3i+i-1}{2}=\dfrac{2+4i}{2}=1+2i$.
\end{exemple}

\methode{Résoudre une équation du second degré à coefficients réels}
Calculer $\Delta$ ; si $\Delta<0$, utiliser $\sqrt{-\Delta}$ et la formule avec $i$.
\begin{exemple}
$z^{2}+z+1=0$ : $\Delta=1-4=-3<0$, donc $z=\dfrac{-1\pm i\sqrt3}{2}$.
\end{exemple}

\methode{Passer de la forme algébrique à la forme trigonométrique}
Calculer $r=|z|$, puis $\theta$ tel que $\cos\theta=a/r$ et $\sin\theta=b/r$ (les \emph{deux}
signes fixent le quadrant).
\begin{exemple}
$z=-1+i\sqrt3$ : $r=\sqrt{1+3}=2$, $\cos\theta=-\frac12$, $\sin\theta=\frac{\sqrt3}2$ donc
$\theta=\frac{2\pi}3$ et $z=2\,e^{\,2i\pi/3}$.
\end{exemple}

\methode{Calculer une puissance avec Moivre}
Mettre sous forme exponentielle $z=re^{i\theta}$, puis $z^{n}=r^{n}e^{in\theta}$ et revenir en
algébrique.
\begin{exemple}
$(1+i)^{8}=\big(\sqrt2\,e^{i\pi/4}\big)^{8}=2^{4}e^{2i\pi}=16$.
\end{exemple}

\methode{Linéariser $\cos^{n}\theta$ ou $\sin^{n}\theta$ avec Euler}
Remplacer $\cos\theta$ ou $\sin\theta$ par leur expression d'Euler, développer (binôme), puis
regrouper les termes conjugués en $\cos(k\theta)$ ou $\sin(k\theta)$.
\begin{exemple}
$\cos^{3}\theta=\Big(\dfrac{e^{i\theta}+e^{-i\theta}}2\Big)^{3}
=\dfrac18\big(e^{3i\theta}+3e^{i\theta}+3e^{-i\theta}+e^{-3i\theta}\big)
=\dfrac14\cos3\theta+\dfrac34\cos\theta$.
\end{exemple}

\methode{Extraire une racine carrée complexe par identification}
Poser $z=x+iy$, écrire $x^{2}-y^{2}=\mathrm{Re}(Z)$, $2xy=\mathrm{Im}(Z)$, $x^{2}+y^{2}=|Z|$ ;
résoudre puis fixer les signes avec $2xy$.
\begin{exemple}
Racines de $-3-4i$ : $|Z|=5$, $x^{2}-y^{2}=-3$, $x^{2}+y^{2}=5$ donc $x^{2}=1$, $y^{2}=4$ ; comme
$2xy=-4<0$, $z=\pm(1-2i)$.
\end{exemple}

\methode{Résoudre $z^{n}=Z$}
Écrire $Z=Re^{i\varphi}$ ; les $n$ racines sont $z_k=R^{1/n}e^{i(\varphi+2k\pi)/n}$,
$k=0,\dots,n-1$.
\begin{exemple}
$z^{3}=8i=8e^{i\pi/2}$ : $z_k=2\,e^{i(\pi/2+2k\pi)/3}$, soit $2e^{i\pi/6},\ 2e^{i5\pi/6},\
2e^{i3\pi/2}=-2i$.
\end{exemple}

\methode{Reconnaître une configuration par un quotient d'affixes}
Calculer $Z=\dfrac{z_C-z_A}{z_B-z_A}$ : son module donne $\dfrac{AC}{AB}$, son argument l'angle
$(\overrightarrow{AB},\overrightarrow{AC})$.
\begin{exemple}
$A(1)$, $B(i)$, $C(-1)$ : $Z=\dfrac{-1-1}{i-1}=1+i$ donc $\frac{AC}{AB}=\sqrt2$ et l'angle vaut
$\frac\pi4$ ; le triangle est rectangle isocèle en $B$.
\end{exemple}

\methode{Déterminer un ensemble de points}
Traduire la condition sur $z$ en distances/arguments : $|z-z_A|=|z-z_B|$ (médiatrice),
$|z-z_A|=r$ (cercle), $\frac{z-z_A}{z-z_B}\in\R$ (droite $(AB)$).
\begin{exemple}
$|z-2|=|z+i|$ : avec $A(2)$, $B(-i)$, c'est $AM=BM$, donc la \textbf{médiatrice} de $[AB]$.
\end{exemple}

% ═══════════════════════════════════════════════════════════════
\section{Exercices d'application}

\rubrique{Forme algébrique, conjugué, module}
\exo{}\diff{1} Écrire sous forme $a+ib$ : $(2-i)(3+i)$, $\dfrac{3+i}{1-i}$ et $\dfrac{1+i}{2-3i}$.

\exo{}\diff{1} Calculer $i^{2027}$, $i^{100}$ et $(1+i)^{2}$. En déduire $(1+i)^{4}$.

\exo{}\diff{1} Déterminer le module de $z=3-4i$, de $z'=-5+12i$ et de $zz'$ (sans développer $zz'$).

\exo{}\diff{2} Soit $z=\dfrac{2+i}{1-2i}$. Montrer que $z$ est un imaginaire pur.

\exo{}\diff{2} Déterminer les réels $x,y$ tels que $(x+iy)(2-i)=3+i$.

\rubrique{Second degré et équations}
\exo{}\diff{1} Résoudre dans $\C$ : $z^{2}-4z+13=0$.

\exo{}\diff{2} Résoudre dans $\C$ : $z^{2}-2z+2=0$, puis donner la forme exponentielle des
solutions.

\exo{}\diff{2} Déterminer les racines carrées du complexe $5-12i$.

\exo{}\diff{3} Résoudre dans $\C$ : $z^{2}-(1+i)z+i=0$ (coefficients complexes).

\rubrique{Formes trigonométrique et exponentielle}
\exo{}\diff{2} Donner la forme exponentielle de $z=-1+i\sqrt3$ et de $z'=-\sqrt3-i$.

\exo{}\diff{2} Calculer $(1+i)^{10}$ à l'aide de la forme exponentielle.

\exo{}\diff{2} Mettre sous forme exponentielle $z=\dfrac{1+i\sqrt3}{1+i}$ et en déduire son module
et un argument.

\exo{}\diff{3} Linéariser $\cos^{3}\theta$ à l'aide des formules d'Euler.

\rubrique{Racines $n$-ièmes}
\exo{}\diff{2} Résoudre $z^{3}=1$ et placer les solutions sur le cercle unité.

\exo{}\diff{2} Résoudre $z^{4}=-16$ dans $\C$ et représenter les solutions.

\exo{}\diff{3} Résoudre $z^{3}=2+2i$ ; donner les solutions sous forme exponentielle.

\rubrique{Interprétation géométrique}
\exo{}\diff{2} Déterminer l'ensemble des points $M$ d'affixe $z$ tels que $|z-2|=|z+i|$.

\exo{}\diff{2} Déterminer l'ensemble des points $M$ tels que $|z-1+2i|=3$.

\exo{}\diff{3} On donne $A,B,C$ d'affixes $z_A=1$, $z_B=i$, $z_C=-1$.
\begin{enumerate}[label=\arabic*)]
\item Calculer $\dfrac{z_C-z_A}{z_B-z_A}$.
\item En déduire la nature du triangle $ABC$.
\end{enumerate}

% ═══════════════════════════════════════════════════════════════
\section{Exercices d'entraînement}

\rubrique{Forme algébrique, conjugué, module}
\exo{}\diff{1} Écrire sous forme algébrique : $(3-2i)(1+4i)$, $\dfrac{5}{2+i}$, $\dfrac{2-i}{3+2i}$.

\exo{}\diff{1} Calculer $i^{2026}$, $i^{-1}$, $i^{2025}+i^{2024}$.

\exo{}\diff{1} Déterminer $\mathrm{Re}(z)$ et $\mathrm{Im}(z)$ pour $z=(2+3i)^{2}$.

\exo{}\diff{2} Montrer que pour tout $z\in\C$, $z+\bar z\in\R$ et $z-\bar z$ est imaginaire pur.

\exo{}\diff{2} Soit $z=\dfrac{1}{1+i}+\dfrac{1}{1-i}$. Montrer que $z$ est réel et le calculer.

\exo{}\diff{2} Déterminer les réels $a,b$ tels que $(a+ib)^{2}=8+6i$.

\exo{}\diff{2} Résoudre dans $\C$ : $\bar z=z^{2}$.

\exo{}\diff{3} Montrer que $|z+z'|^{2}+|z-z'|^{2}=2\big(|z|^{2}+|z'|^{2}\big)$ (identité du
parallélogramme).

\rubrique{Second degré et équations}
\exo{}\diff{1} Résoudre : $z^{2}+9=0$ ; $z^{2}+2z+5=0$.

\exo{}\diff{2} Résoudre dans $\C$ : $z^{2}-6z+25=0$, puis $2z^{2}-2z+1=0$.

\exo{}\diff{2} Déterminer les racines carrées de $-7+24i$.

\exo{}\diff{2} Résoudre $z^{2}+(2-i)z+(3-i)=0$.

\exo{}\diff{3} Résoudre dans $\C$ : $z^{3}-1=0$, puis en déduire la factorisation de $z^{3}-8$.

\exo{}\diff{3} Soit $P(z)=z^{3}-3z^{2}+4z-2$. Vérifier que $z=1$ est racine, factoriser $P$ et
résoudre $P(z)=0$ dans $\C$.

\exo{}\diff{3} Résoudre dans $\C$ : $z^{4}+z^{2}+1=0$ (poser $Z=z^{2}$).

\rubrique{Formes trigonométrique et exponentielle}
\exo{}\diff{1} Donner la forme exponentielle de $1+i$, de $-i$, de $-3$ et de $1-i\sqrt3$.

\exo{}\diff{2} Calculer $(1-i\sqrt3)^{6}$.

\exo{}\diff{2} Mettre $\dfrac{(1+i)^{5}}{(1-i)^{3}}$ sous forme algébrique.

\exo{}\diff{2} Déterminer le module et un argument de $z=(\sqrt3+i)(1+i)$.

\exo{}\diff{3} Linéariser $\sin^{2}\theta\cos^{2}\theta$, puis $\sin^{3}\theta$.

\exo{}\diff{3} À l'aide de Moivre, exprimer $\cos3\theta$ et $\sin3\theta$ en fonction de
$\cos\theta$ et $\sin\theta$.

\exo{}\diff{3} Calculer la somme $S=\sum_{k=0}^{n-1}\cos(k\theta)$ pour $\theta\neq0\ [2\pi]$ (somme
de série géométrique de raison $e^{i\theta}$).

\rubrique{Racines $n$-ièmes}
\exo{}\diff{2} Résoudre $z^{4}=1$ et $z^{6}=1$ ; représenter les solutions.

\exo{}\diff{2} Résoudre $z^{3}=-27$.

\exo{}\diff{3} Résoudre $z^{4}=-8+8i\sqrt3$.

\exo{}\diff{3} Soit $j=e^{2i\pi/3}$. Montrer que $1+j+j^{2}=0$ et que $j^{3}=1$.

\rubrique{Interprétation géométrique}
\exo{}\diff{2} Déterminer l'ensemble des $M(z)$ tels que $|z-i|=|z+1|$.

\exo{}\diff{2} Déterminer l'ensemble des $M(z)$ tels que $|2z-1|=4$.

\exo{}\diff{2} Soient $A(2+i)$, $B(-1+3i)$. Calculer $AB$ et l'affixe du milieu de $[AB]$.

\exo{}\diff{3} Soient $A(1+i)$, $B(3-i)$, $C(2+4i)$. Le triangle $ABC$ est-il rectangle ? isocèle ?

\exo{}\diff{3} Déterminer l'ensemble des $M(z)$ tels que $\dfrac{z-1}{z+1}$ soit imaginaire pur.

\exo{}\diff{3} Déterminer l'ensemble des $M(z)$, $z\neq2$, tels que $\dfrac{z-i}{z-2}$ soit réel.

% ═══════════════════════════════════════════════════════════════
\section{Sujets type Baccalauréat}

\sujetbac[Plan complexe, second degré, configuration]
On considère le polynôme $P(z)=z^{3}-(4+i)z^{2}+(7+i)z-4$, et l'on note $A$, $B$, $C$ les points
d'affixes les racines de $P$.
\begin{enumerate}[label=\arabic*)]
\item Vérifier que $z_0=1$ est une racine de $P$.
\item Déterminer les réels $a,b$ tels que $P(z)=(z-1)\big(z^{2}+az+b\big)$ avec $a,b\in\C$.
\item Résoudre $P(z)=0$ dans $\C$.
\item On pose $A(1)$, $B(2+i)$, $C(1+i)$ avec les racines trouvées. Calculer
$\dfrac{z_C-z_A}{z_B-z_A}$ et en déduire la nature du triangle $ABC$.
\end{enumerate}

\sujetbac[Forme exponentielle, Moivre, lieu de points]
Soit $z=1+i\sqrt3$.
\begin{enumerate}[label=\arabic*)]
\item Donner le module et un argument de $z$, puis sa forme exponentielle.
\item En déduire la forme exponentielle de $z^{6}$ et $z^{2026}$.
\item Résoudre dans $\C$ l'équation $Z^{2}=z$, en donnant les solutions sous forme exponentielle.
\item Déterminer l'ensemble des points $M$ d'affixe $w$ tels que $|w-z|=|z|$.
\end{enumerate}

\sujetbac[Racines de l'unité et somme]
Soit $n\ge2$ un entier et $\omega_k=e^{2ik\pi/n}$, $k=0,\dots,n-1$, les racines $n$-ièmes de
l'unité.
\begin{enumerate}[label=\arabic*)]
\item Montrer que $\displaystyle\sum_{k=0}^{n-1}\omega_k=0$.
\item Application $n=3$ : on pose $j=e^{2i\pi/3}$. Calculer $1+j+j^{2}$ et $j^{2}+\bar j$.
\item Montrer que les images de $1,j,j^{2}$ forment un triangle équilatéral et préciser son centre
de gravité.
\end{enumerate}

% ═══════════════════════════════════════════════════════════════
\section{Corrigés}

\corrige{de l'exercice 1}
$(2-i)(3+i)=6+2i-3i-i^{2}=7-i$.\;
$\dfrac{3+i}{1-i}=\dfrac{(3+i)(1+i)}{2}=\dfrac{2+4i}{2}=1+2i$.\;
$\dfrac{1+i}{2-3i}=\dfrac{(1+i)(2+3i)}{13}=\dfrac{2+3i+2i-3}{13}=\dfrac{-1+5i}{13}
=-\dfrac1{13}+\dfrac{5}{13}i$.

\corrige{de l'exercice 2}
$2027=4\times506+3$ donc $i^{2027}=i^{3}=-i$.\; $100=4\times25$ donc $i^{100}=1$.\;
$(1+i)^{2}=1+2i-1=2i$, d'où $(1+i)^{4}=(2i)^{2}=-4$.

\corrige{de l'exercice 3}
$|z|=\sqrt{9+16}=5$,\; $|z'|=\sqrt{25+144}=13$,\; donc $|zz'|=|z|\,|z'|=65$.

\corrige{de l'exercice 4}
$z=\dfrac{2+i}{1-2i}=\dfrac{(2+i)(1+2i)}{5}=\dfrac{2+4i+i-2}{5}=\dfrac{5i}{5}=i$ : c'est un
imaginaire pur.

\corrige{de l'exercice 5}
$(x+iy)(2-i)=(2x+y)+i(2y-x)=3+i$. Par identification : $2x+y=3$ et $-x+2y=1$. On résout :
$y=3-2x$, donc $-x+6-4x=1$, soit $5x=5$, $x=1$ et $y=1$.

\corrige{de l'exercice 6}
$\Delta=16-52=-36<0$, $\sqrt{-\Delta}=6$ : $z=\dfrac{4\pm6i}{2}=2\pm3i$.

\corrige{de l'exercice 7}
$\Delta=4-8=-4<0$, $\sqrt{-\Delta}=2$ : $z=\dfrac{2\pm2i}{2}=1\pm i$. Or $1+i=\sqrt2\,e^{i\pi/4}$ et
$1-i=\sqrt2\,e^{-i\pi/4}$.

\corrige{de l'exercice 8}
On cherche $z=x+iy$ avec $z^{2}=5-12i$. Alors $x^{2}-y^{2}=5$, $2xy=-12$, $x^{2}+y^{2}=|5-12i|=13$.
D'où $x^{2}=9$, $y^{2}=4$ ; comme $2xy=-12<0$, $x$ et $y$ sont de signes opposés :
$z=\pm(3-2i)$.

\corrige{de l'exercice 9}
$\Delta=(1+i)^{2}-4i=2i-4i=-2i$. Une racine carrée $\delta=x+iy$ de $-2i$ : $x^{2}-y^{2}=0$,
$2xy=-2$, $x^{2}+y^{2}=2$, d'où $x^{2}=y^{2}=1$ et $xy<0$ : $\delta=1-i$. Alors
$z=\dfrac{(1+i)\pm(1-i)}{2}$, soit $z=1$ ou $z=i$.

\corrige{de l'exercice 10}
$z=-1+i\sqrt3$ : $r=\sqrt{1+3}=2$, $\cos\theta=-\frac12$, $\sin\theta=\frac{\sqrt3}2$, donc
$\theta=\frac{2\pi}3$ et $z=2e^{2i\pi/3}$.\;
$z'=-\sqrt3-i$ : $r=\sqrt{3+1}=2$, $\cos\theta=-\frac{\sqrt3}2$, $\sin\theta=-\frac12$, donc
$\theta=-\frac{5\pi}6$ et $z'=2e^{-5i\pi/6}$.

\corrige{de l'exercice 11}
$1+i=\sqrt2\,e^{i\pi/4}$, donc $(1+i)^{10}=2^{5}e^{10i\pi/4}=32\,e^{5i\pi/2}=32\,e^{i\pi/2}=32i$.

\corrige{de l'exercice 12}
$1+i\sqrt3=2e^{i\pi/3}$ et $1+i=\sqrt2\,e^{i\pi/4}$, donc
$z=\dfrac{2}{\sqrt2}\,e^{i(\pi/3-\pi/4)}=\sqrt2\,e^{i\pi/12}$. Module $\sqrt2$, argument
$\frac\pi{12}$.

\corrige{de l'exercice 13}
$\cos\theta=\dfrac{e^{i\theta}+e^{-i\theta}}2$, donc
$\cos^{3}\theta=\dfrac18\big(e^{i\theta}+e^{-i\theta}\big)^{3}
=\dfrac18\big(e^{3i\theta}+3e^{i\theta}+3e^{-i\theta}+e^{-3i\theta}\big)
=\dfrac14\cos3\theta+\dfrac34\cos\theta$.

\corrige{de l'exercice 14}
$z^{3}=1$ : $z_k=e^{2ik\pi/3}$, $k=0,1,2$, soit
$1,\ -\frac12+i\frac{\sqrt3}2,\ -\frac12-i\frac{\sqrt3}2$ : trois sommets d'un triangle équilatéral
sur le cercle unité.

\corrige{de l'exercice 15}
$-16=16e^{i\pi}$, donc $z_k=16^{1/4}e^{i(\pi+2k\pi)/4}=2e^{i(\pi/4+k\pi/2)}$, $k=0,1,2,3$ :
\[
2e^{i\pi/4}=\sqrt2(1+i),\ 2e^{i3\pi/4}=\sqrt2(-1+i),\ 2e^{i5\pi/4}=\sqrt2(-1-i),\
2e^{i7\pi/4}=\sqrt2(1-i),
\]
sommets d'un carré centré en $O$.

\corrige{de l'exercice 16}
$2+2i=2\sqrt2\,e^{i\pi/4}$, donc
$z_k=(2\sqrt2)^{1/3}e^{i(\pi/4+2k\pi)/3}=\sqrt2\,e^{i(\pi/12+2k\pi/3)}$, $k=0,1,2$ (car
$(2\sqrt2)^{1/3}=(2^{3/2})^{1/3}=2^{1/2}=\sqrt2$). Soit
$\sqrt2\,e^{i\pi/12},\ \sqrt2\,e^{3i\pi/4},\ \sqrt2\,e^{i17\pi/12}$.

\corrige{de l'exercice 17}
$|z-2|=|z+i|$ s'écrit $AM=BM$ avec $A(2)$ et $B(-i)$ : c'est la \textbf{médiatrice} du segment
$[AB]$.

\corrige{de l'exercice 18}
$|z-(1-2i)|=3$ : cercle de centre $\Omega(1-2i)$ et de rayon $3$.

\corrige{de l'exercice 19}
1) $\dfrac{z_C-z_A}{z_B-z_A}=\dfrac{-1-1}{i-1}=\dfrac{-2}{i-1}=\dfrac{-2(-1-i)}{|i-1|^{2}}
=\dfrac{2+2i}{2}=1+i$.\;
2) $|1+i|=\sqrt2$ donc $\dfrac{AC}{AB}=\sqrt2$, et $\arg(1+i)=\frac\pi4$ donc
$(\overrightarrow{AB},\overrightarrow{AC})=\frac\pi4$. On vérifie de plus $AB=|i-1|=\sqrt2$ et
$BC=|{-1}-i|=\sqrt2$ : le triangle $ABC$ est \textbf{rectangle isocèle en $B$}.

\corrige{du sujet 1}
1) $P(1)=1-(4+i)+(7+i)-4=1-4-i+7+i-4=0$ : $z_0=1$ est racine.\;
2) Division par $(z-1)$ : $P(z)=(z-1)\big(z^{2}-(3+i)z+4\big)$, donc $a=-(3+i)$, $b=4$.\;
3) Reste à résoudre $z^{2}-(3+i)z+4=0$ : $\Delta=(3+i)^{2}-16=9+6i-1-16=-8+6i$. Racine carrée
$\delta=x+iy$ : $x^{2}-y^{2}=-8$, $2xy=6$, $x^{2}+y^{2}=|{-8+6i}|=10$, d'où $x^{2}=1$, $y^{2}=9$ et
$xy>0$ : $\delta=1+3i$. Alors $z=\dfrac{(3+i)\pm(1+3i)}{2}$, soit $z=2+2i$ ou $z=1-i$. Les racines
de $P$ sont donc $1,\ 2+2i,\ 1-i$.\;
4) Avec $A(1)$, $B(2+2i)$, $C(1-i)$ :
$\dfrac{z_C-z_A}{z_B-z_A}=\dfrac{-i}{1+2i}=\dfrac{-i(1-2i)}{5}=\dfrac{-i-2}{5}=\dfrac{-2-i}{5}$.
Ce nombre n'est ni réel ni imaginaire pur : le triangle est quelconque. \emph{(Avec l'énoncé
proposé $B(2+i)$, $C(1+i)$ on aurait $\frac{z_C-z_A}{z_B-z_A}=\frac{i}{1+i}=\frac{i(1-i)}{2}
=\frac{1+i}{2}$ : module $\frac{\sqrt2}2$, argument $\frac\pi4$ — triangle rectangle isocèle
direct en $A$.)}

\corrige{du sujet 2}
1) $z=1+i\sqrt3$ : $|z|=\sqrt{1+3}=2$, $\cos\theta=\frac12$, $\sin\theta=\frac{\sqrt3}2$ donc
$\theta=\frac\pi3$ et $z=2e^{i\pi/3}$.\;
2) $z^{6}=2^{6}e^{6i\pi/3}=64\,e^{2i\pi}=64$.\; $z^{2026}=2^{2026}e^{2026i\pi/3}$ ; or
$2026=3\times675+1$ donc $e^{2026i\pi/3}=e^{i\pi/3}$ et $z^{2026}=2^{2026}e^{i\pi/3}$.\;
3) $Z^{2}=2e^{i\pi/3}$ : $Z_k=\sqrt2\,e^{i(\pi/3+2k\pi)/2}=\sqrt2\,e^{i(\pi/6+k\pi)}$, $k=0,1$, soit
$Z=\sqrt2\,e^{i\pi/6}$ et $Z=\sqrt2\,e^{i7\pi/6}=-\sqrt2\,e^{i\pi/6}$.\;
4) $|w-z|=|z|=2$ : cercle de centre $\Omega$ d'affixe $z=1+i\sqrt3$ et de rayon $2$ (il passe par
$O$ car $|0-z|=|z|=2$).

\corrige{du sujet 3}
1) Pour $\omega=e^{2i\pi/n}\neq1$, $\displaystyle\sum_{k=0}^{n-1}\omega^{k}
=\dfrac{\omega^{n}-1}{\omega-1}=\dfrac{1-1}{\omega-1}=0$ (somme géométrique, $\omega^{n}=1$).\;
2) $1+j+j^{2}=0$ (cas $n=3$). Comme $|j|=1$, $\bar j=\frac1j=j^{2}$ ; donc
$j^{2}+\bar j=j^{2}+j^{2}=2j^{2}$. \emph{(On peut aussi écrire $j^{2}+\bar j=j^{2}+j^{2}=2j^{2}
=-2-2j$ d'après $1+j+j^{2}=0$.)}\;
3) Les points $1,j,j^{2}$ sont sur le cercle unité, séparés de $\frac{2\pi}3$ : c'est un triangle
\textbf{équilatéral}. Son centre de gravité a pour affixe $\frac{1+j+j^{2}}3=0$ : c'est l'origine
$O$ (centre du cercle).
